\section{Methodology}
\label{sec:method}

\subsection{Additive prompt levels}
\label{sec:levels}

For each task we define seven prompt templates $T_1 \subset T_2 \subset
\cdots \subset T_7$, where $\subset$ denotes strict token superset.
Each level appends one layer of elaboration to all previous content;
no layer modifies earlier text.
Table~\ref{tab:levels} shows the layer structure for all six tasks.

\begin{table}[ht]
\centering
\caption{Additive prompt layers by task. Each row adds content to all previous rows.}
\label{tab:levels}
\small
\resizebox{\textwidth}{!}{
\begin{tabular}{cllllllll}
\toprule
Lv & Layer added & Classification & Prod.\ Extraction & QA & Math Reasoning & Summarization & Instr.\ Following \\
\midrule
1 & Bare input   & Raw text only & Product text & Question only & Problem statement & Article text & Topic prompt \\
2 & Task label   & ``Classify sentiment'' & List field names & ``Answer the question'' & ``Solve the problem'' & ``Summarize this'' & Task + constraints \\
3 & Format spec  & ``Pos/Neg/Neu'' & JSON format & ``Answer concisely'' & ``Give a number'' & Length target & Format specification \\
4 & Definitions  & Label definitions & Field definitions & Domain context & Show-work instr. & Content guidance & Constraint definitions \\
5 & Persona      & Expert analyst & Expert extractor & Domain expert & Math tutor & Journalist & Writing expert \\
6 & Guidelines   & Mixed/ironic rules & Edge cases & Uncertainty & Step-by-step rules & Quality rules & Formatting rules \\
7 & Example      & Worked example & Worked example & Worked example & Worked example & Worked example & Worked example \\
\bottomrule
\end{tabular}
}
\end{table}

The additive design is critical: because level $k+1$ contains all tokens of
level $k$ plus exactly one new layer, any quality difference between adjacent
levels is attributable solely to that layer.
Importantly, no layer is filler; every added layer is meaningful instruction
that could plausibly improve performance.

\subsection{Tasks}
\label{sec:tasks}

We study six tasks spanning a structured-to-open spectrum:

\begin{itemize}
  \item \textbf{Classification}: sentiment classification (positive / negative / neutral)
    over 20 product reviews and social media posts.
    Heuristic evaluator: correct label present in response.
  \item \textbf{Product extraction}: extract name, price, brand, category from
    product listings in JSON format.
    20 examples with varying difficulty (buried prices, implicit brand names).
    Heuristic: fraction of four fields correctly matched.
  \item \textbf{Question answering (QA)}: factoid questions across geography,
    science, and history. 20 examples with ground-truth strings.
    Heuristic: word-overlap recall (ROUGE-1).
  \item \textbf{Math reasoning}: 2--3 step arithmetic, rate, and geometry
    problems. 20 examples (10 easy, 10 medium difficulty).
    Heuristic: exact numerical match on the last number in the response.
  \item \textbf{Summarization}: abstractive summarisation of news passages.
    20 examples with reference summaries.
    Heuristic: length score + ROUGE-1 recall.
  \item \textbf{Instruction following}: responses must satisfy explicit
    structural constraints (bullet points, numbered list, single word).
    20 examples, each with 1--2 constraints.
    Heuristic: fraction of constraints satisfied.
\end{itemize}

Tasks are grouped \textit{post hoc} into \textbf{structured} (classification,
product extraction) and \textbf{open-ended} (QA, math, summarization,
instruction following) based on the analysis in Section~\ref{sec:results}.

\subsection{Models}
\label{sec:models}

We evaluate seven LLMs across four providers, covering a broad capability and
size range (Table~\ref{tab:models}).
All models are accessed via production APIs with default temperature settings.
Each response is capped at 512 output tokens.

\begin{table}[t]
\centering
\caption{Models evaluated. Parameter counts are listed where publicly
  disclosed; ``undiscl.''\ indicates the provider has not released this
  information.}
\label{tab:models}
\small
\begin{tabular}{lllr}
\toprule
Model (alias) & Provider & API & Parameters \\
\midrule
llama-3.1-8b-instant    & Meta       & Groq      & 8B        \\
qwen/qwen3-32b          & Alibaba    & Groq      & 32B       \\
llama-3.3-70b-versatile & Meta       & Groq      & 70B       \\
gemini-2.0-flash        & Google     & Gemini    & undiscl.  \\
moonshotai/kimi-k2      & Moonshot   & Groq      & undiscl.  \\
gpt-4o-mini             & OpenAI     & OpenAI    & undiscl.  \\
claude-haiku-4-5        & Anthropic  & Anthropic & undiscl.  \\
\bottomrule
\end{tabular}
\end{table}

\subsection{Evaluation}
\label{sec:eval}

\paragraph{LLM judge.}
Following \citet{zheng2023judging}, who show that strong LLM judges match
human preferences at over 80\% agreement, all 5,875 responses are scored by
\texttt{gpt-4o-mini} on four dimensions (correctness, completeness,
reasoning, conciseness), each rated 1--5.
The aggregate quality score is the normalised sum:
$q = \frac{\sum_i s_i}{20} \in [0, 1]$.
We provide task-specific rubrics to the judge (e.g., for classification: focus
on whether the label matches; for product extraction: whether all four fields
are present and correct).

\paragraph{Heuristic validation.}
We additionally validate the LLM judge against a task-specific heuristic
evaluator (described in Section~\ref{sec:tasks}) on the full dataset and find
Pearson $r = 0.986$, confirming that judge scores are reliable and not a
source of noise in the results.

\paragraph{Aggregation.}
For each (model, task, level) triplet, we aggregate quality across all 20
examples by taking the mean.
This yields 7 data points per (model, task) curve.

\subsection{Curve fitting}
\label{sec:fitting}

For each (model, task) pair we fit two functional forms to the 7 level means:
\begin{equation}
  q_{\log}(x) = a \log(bx) + c, \qquad
  q_{\text{sig}}(x) = \frac{L}{1 + e^{-k(x - x_0)}} + c
  \label{eq:curves}
\end{equation}
where $x$ is the mean prompt token count at each level.
We select the form with lower RMSE.
The \textbf{saturation point} $T^*$ is defined as the smallest $x$ such that:
\begin{equation}
  q(x) \geq 0.95 \times \max_{x} q(x)
  \label{eq:satpoint}
\end{equation}

\paragraph{Null model F-test.}
To test whether a fitted curve is significantly better than a flat line
($q(x) = \bar{q}$), we compute:
\begin{equation}
  F = \frac{(SS_{\text{null}} - SS_{\text{fit}}) / (k - 1)}{SS_{\text{fit}} / (n - k)}
  \label{eq:ftest}
\end{equation}
where $k \in \{3, 4\}$ is the number of curve parameters (log or sigmoid),
$n = 7$ is the number of level means, and we use $\alpha = 0.05$.
The null model is a constant $q(x) = \bar{q}$, \ie, prompt length has no
effect on quality; a significant F-test rejects this in favour of the fitted
curve.

\paragraph{Bootstrap confidence intervals.}
We estimate 95\% CIs on $T^*$ via 1,000-iteration bootstrap.
On each iteration we resample the 20 example indices \emph{with replacement}
and apply the same resample across all seven levels, preserving the paired
structure.
We re-aggregate, re-fit, and re-compute $T^*$ for each bootstrap replicate;
the 2.5th and 97.5th percentiles of the resulting distribution form the CI.

\subsection{Experimental scale}
\label{sec:scale}

Total API calls: $7 \times 6 \times 7 \times 20 = 5{,}880$, with 5,875
successful records (99.9\%).
All prompts, code, and data are included in the supplementary material.
